
%%% Local Variables:
%%% mode: latex
%%% TeX-master: "../main"
%%% End:

\begin{ack}
  转瞬间我的研究生生涯即将结束。遥想三年前的保研过程中机缘巧合下来到了国防科技大学这一可以说是全中国最特殊的的大学就读，
  我从一开始的不习惯到逐渐适应，再到现在的依依不舍，国防科技大学的这段求学经历无疑是我人生中最特殊的宝贵回忆。
  一路走来，有过迷茫，有过挫折，但更多的是收获和成长。回首这段经历，我想对所有给予我帮助和支持的人表达最诚挚的感谢。

  首先衷心感谢我的导师占荣辉老师对本人的精心指导。占老师事必躬亲，踏实严谨又为人和善，
  在缘分下遇到了占老师是我研究生阶段最大的幸运。
  占老师不仅在学术上给予我悉心指导，在其他事务上也绝不麻烦学生，宁愿自己加加班，杂事也都尽量会自己处理。
  占老师的严谨治学态度和对学生的关爱让我受益匪浅，以及这种踏实做事、平和待人的品质，我会永远铭记在心。
  也要感谢教研室的张军老师、王威老师、刘盛启老师、回丙伟老师、欧建平老师、
  吴中杰老师以及陈诗琪老师在硕士期间对我学习和生活上的帮助和教导。

  感谢师兄师姐们的无私帮助和支持。感谢师兄师姐们在学术上给予的指导和建议，在生活上给予的关心和帮助。
  尤其感谢郭越师兄、张会强大师兄和赵久超大大师兄带领我进入ATR1室，带我熟悉环境，帮助我成长。
  在我研究生阶段的提供了诸多学术指导和生活帮助。也感谢各位同门以及师弟师妹们的陪伴和支持，
  作为一起同甘共苦的小伙伴，感谢你们在学术上给予的建议，在繁琐的各项事务上给予的帮助。

  感谢我的三位舍友刘振威、潘永雷和何锦松的陪伴和支持。在国防科大这样的特殊环境中，我们抱团取暖，
  互帮互助。我靠着和你们一起释放压力，分享生活，才能逐渐找到自己的舒适区。
  最后，感谢我的家人对我的支持和鼓励。感谢父母一直以来的理解、支持和信任，
  感谢他们在我求学路上的无私奉献和鼓励。

  人生就是在不断找寻意义的旅途，从一个个阶段的意义中不断前行，找寻自己真正的人生意义。
  硕士阶段的目标已经结束，我将继续前行，去追寻下一个阶段的我。

\end{ack}
